%
% halbkreis-ohne-pol.tex -- Halbkreiskontur ohne Polstelle auf der Achse
%
% (c) 2026 Damien Flury, OST Ostschweizer Fachhochschule
%
\documentclass[tikz]{standalone}
\usepackage{amsmath}
\usepackage{times}
\usepackage{txfonts}
\usetikzlibrary{arrows,math,decorations.markings}
\definecolor{darkblue}{rgb}{0,0,0.7}
\definecolor{darkred}{rgb}{0.7,0,0}
\begin{document}
\def\skala{1}
\begin{tikzpicture}[>=latex,thick,scale=\skala,
	midarrow/.style={
		postaction={decorate},
		decoration={markings,mark=at position #1 with {\arrow[line width=1.2pt]{latex}}}}]

% Parameter der Zeichnung
\def\R{3.2}
\def\i{1.3}

% Achsen
\draw[->] (-\R-0.6,0) -- (\R+0.7,0) coordinate[label={$\operatorname{Re}z$}];
\draw[->] (0,-\i-0.6) -- (0,\R+0.7) coordinate[label={right:$\operatorname{Im}z$}];

% grosser Halbkreis C_R, im positiven Sinn durchlaufen
\draw[color=darkblue,line width=1.4pt,midarrow=0.5,midarrow=0.85]
	(\R,0) arc[start angle=0, end angle=180, radius=\R];
\node[color=darkblue] at (135:\R+0.42) {$C_R$};

% Strecke von -R bis R auf der reellen Achse
\draw[color=darkblue,line width=1.4pt,midarrow=0.3,midarrow=0.75]
	(-\R,0) -- (\R,0);

% Pol i in der oberen Halbebene, vom Weg umlaufen
\fill[color=darkred] (0,\i) circle[radius=0.06];
\node[color=darkred] at (0,\i) [above right] {$i$};

% Pol -i in der unteren Halbebene, nicht umlaufen
\draw[color=gray,line width=0.8pt] (0,-\i) circle[radius=0.06];
\node[color=gray] at (0,-\i) [right=3pt] {$-i$};

% Beschriftung der Integrationsgrenzen
\draw (-\R,-0.08) -- (-\R,0.08);
\node at (-\R,-0.42) {$-R$};
\draw (\R,-0.08) -- (\R,0.08);
\node at (\R,-0.42) {$R$};

\end{tikzpicture}
\end{document}
